\section{Điều khiển, di chuyển và xoay khối}

\subsection{Cách chơi và hệ thống phím}

Mã nguồn được tổ chức theo kiến trúc mô-đun trong thư mục \texttt{src}.
\texttt{Board} chịu trách nhiệm lưu bàn chơi, kiểm tra va chạm, khóa khối và
xóa dòng. \texttt{Game} điều phối trạng thái khối đang rơi và các thao tác
điều khiển. \texttt{Input}, \texttt{Renderer} và \texttt{Score} lần lượt xử lý
bàn phím, hiển thị console và tính điểm.

Hàm kiểm tra va chạm chỉ cho phép khối được đặt trong vùng chơi và không
chồng lên các ô đã khóa. Phép xoay theo chiều kim đồng hồ được cài đặt bằng
công thức:
\[
    \mathrm{new}[i][j] = \mathrm{current}[3-j][i].
\]

Người chơi dùng phím \texttt{A}/\texttt{D} để di chuyển trái/phải,
\texttt{X} để rơi nhanh, \texttt{W} để xoay và \texttt{Q} để thoát. Khi phép
xoay bị chặn cạnh tường hoặc khối đã khóa, chương trình thử các độ lệch ngang
nhỏ trong khoảng từ $-2$ đến $2$. Nếu không có vị trí hợp lệ, khối giữ nguyên
trạng thái hiện tại để tránh xuyên tường. Các vị trí thử nghiệm gồm cả độ lệch
ngang và dọc nhỏ, giúp xoay được khi khối đang sát sàn hoặc một cấu trúc đã
khóa.

Renderer hiển thị ghost piece tại vị trí thấp nhất mà khối có thể hạ xuống.
Sau mỗi 10 dòng bị xóa, level tăng và thời gian rơi tự động giảm, nhưng
không thấp hơn 80 mili giây.
